% Generated draft: requires mathematical and source review.
\documentclass[11pt]{article}
\usepackage[margin=25mm]{geometry}
\usepackage{fontspec}
\setmainfont{Latin Modern Roman}
\usepackage{amsmath,amssymb,mathtools,unicode-math}
\setmathfont{Latin Modern Math}
\usepackage{enumitem,needspace,xurl,hyperref}
\setlength{\parindent}{0pt}
\setlength{\parskip}{5pt}
\setlength{\emergencystretch}{3em}
\newcommand{\N}{\mathbb{N}}
\newcommand{\Z}{\mathbb{Z}}
\newcommand{\Q}{\mathbb{Q}}
\newcommand{\R}{\mathbb{R}}
\newcommand{\C}{\mathbb{C}}
\newcommand{\F}{\mathbb{F}}
\newcommand{\A}{\mathbb{A}}
\newcommand{\PP}{\mathbb{P}}
\newcommand{\E}{\mathbb{E}}
\newcommand{\eps}{\varepsilon}
\newcommand{\dd}{\,\mathrm{d}}
\begin{document}
\section*{UnsolvedMath ID: 800359}
\textbf{Source problem code:} catalogue entry
\subsection{Definitions and mathematical statement}
Let \(T_n\) denote the region and let ``stones'' denote the allowed tiles intended in the supplied source. The supplied record does not define either object, the permitted tile orientations, or the equivalence convention for counting tilings. Consequently, the following formulation is conditional on recovering those definitions from the cited source; it does not specify a complete tiling model.

With those definitions fixed, let
\[
a_n=\#\{\text{tilings of }T_n\text{ by stones}\},
\qquad
S=\{n\in\N:n\equiv 0,2,9,11\pmod{12}\}.
\]
For a positive integer \(a\), define
\[
v_2(a)=\max\{k\in\Z_{\geq 0}:2^k\mid a\},
\]
and set \(v_2(0)=+\infty\). The source asks for the numbers \(a_n\), and specifically asks whether
\[
\lim_{\substack{n\to\infty\\n\in S}}a_n=0
\quad\text{in the }2\text{-adic topology}.
\]
Equivalently, is it true that
\[
\forall m\in\Z_{\geq 1}\;\exists N\in\N\;\forall n\in S,
\qquad n\geq N\Longrightarrow 2^m\mid a_n?
\]
This is also equivalent to \(v_2(a_n)\to+\infty\) along \(S\), with the convention above for zero counts.

As motivation, the supplied source associates the counts with OEIS A334875 and reports that the valuations of its nonzero terms begin
\[
0,0,1,3,2,3,4,3,4,3,5,8,6,8,\ldots.
\]
The record does not supply the indices corresponding to these values. Its comparison with the decreasing divisibility frequencies of randomly chosen positive integers is motivational, not a probabilistic hypothesis. The registry labels the problem open; that label is supplied source data, not an independently verified status assessment.
\subsection{Short English statement}
Count the tilings of the intended triangular-grid region \(T_n\) by the intended stone tiles. Must these counts eventually be divisible by every fixed power of \(2\), as \(n\) tends to infinity through the residue classes \(0,2,9,11\pmod{12}\)? The supplied record omits the definitions needed to determine the tiling counts.
\subsection{Sources}
\begin{itemize}
\item Local UnsolvedMath registry, permanent ID \(800359\), source problem code \texttt{catalogue entry}, title ``Trimer Covers in the Triangular Grid --- Problem 20.'' The supplied record attributes the problem to Propp (2022) and labels it open.
\item Propp, \emph{Trimer Covers in the Triangular Grid: Twenty Mostly Open Problems} (2022), source item recorded as ``Problem 20 (\#11).'' Supplied URL: \url{https://arxiv.org/abs/2206.06472}. The background records access on 2026-10-16 and extraction by \texttt{arxiv\_latex}; these are supplied provenance statements, not claims of independent access or verification.
\item OEIS entry A334875, \url{https://oeis.org/A334875}, cited in the supplied statement for counts at small values of \(n\). Its contents and indexing have not been independently verified here.
\end{itemize}
\end{document}
